\newpage

% Elemento opcional. Consiste na relação alfabética das abreviaturas e siglas utilizadas no texto, seguidas das palavras ou expressões correspondentes grafadas por extenso separadas por um traço. Recomenda-se a elaboração de lista própria para cada tipo. 

% Exemplo: ABNT – Associação Brasileira de Normas Técnicas, IBGE – Instituto Brasileiro de Geografia e Estatística, RCC – Resíduo da Construção Civil 

\printglossary[type=\acronymtype, title={Lista de Abreviaturas e Siglas}]

% No texto, as abreviaturas que aparecem pela primeira vez, em ordem de leitura, devem ser escritas por extenso, seguida então da abreviatura entre parênteses. Se a abreviatura se repetir no texto, poderá então ser utilizada sem a escrita por extenso. 

% Exemplo: A Associação Brasileira de Normas Técnicas (ABNT) é responsável pelas publicações das normas técnicas regulamentadoras do país.

% Lista de Abreviaturas e Siglas
\newacronym{acp}{ACP}{Análise de Componentes Principais}
\newacronym{rna}{RNA}{Redes Neurais Artificiais}
\newacronym{teo}{OTE}{Ótimo técnico-econômico}

% Para usar no documento: \gls{acp} 